\refstepcounter{section}
\section*{Lecture 23: Continuous Symmetry-II}
\addcontentsline{toc}{section}{Lecture 23: Continuous Symmetry-II}
Let us understand the gapped and gapless modes more closely. For that, consider the fluctuation $\boldsymbol{\psi}$ to the ground state as defined before such that we can write, 
$$\bphi  = \mqty[m\tund{0}+\psi_1\\ \psi_2] $$ Instead of writing the fluctuation in terms of the transverse component $\psi_2(\vb{x})$, we write it as a local change in the direction.
\begin{figure}[H]
    \centering
    \input{pics/rot.tex}
\end{figure}
\noindent
We can define $\tan \theta = \frac{\psi_2}{m\tund{0}+\psi_1}$. Considering the fluctuations to be small, we can effectively take $\tan \theta \approx \theta$ and $\abs{\psi_1}\ll \abs{m\tund{0}}$ which gives us $\psi_2 \simeq m\tund{0}\theta$. Then the transverse fluctuation is completely characterised by $\theta$. From Eq. \eqref{landaucontsymm}, we can use the same \textit{equipartition theorem} to find that, $$V\SG_{\perp}(k) = \expval{\abs{\widetilde{\theta}_{\vb{k}}}^2} = \frac{V\kb T}{\gamma m\tund{0}^2 k^2} \ \longrightarrow \ \SG_{\perp}(k) = \frac{\kb T}{\gamma m\tund{0}^2 k^2}$$
We now calculate the correlation function in real space, which gives us
\begin{equation}
    \SG_{\perp}(\vb{r}) = \int \frac{\dd^d k}{(2\pi)^d}\ \frac{\kb T}{\gamma m\tund{0}^2 k^2} e^{i\vb{k}\cdot \vb{r}} = - \frac{\kb T}{\gamma m\tund{0}^2} C_d
\end{equation}
where we defined this integral in terms of $C_d$ which is the $d$-dimensional Coulomb potential for a unit charge at the origin, which can be easily seen, since it satisfies the Poisson equation,
\begin{equation}
    \nabla^2 C_d(\vb{r}) = - \int  \frac{\dd^d k}{(2\pi)^d}\ \frac{1}{ \cancel{k^2}}\cdot (-\cancel{k^2})\cdot  e^{i\vb{k}\cdot \vb{r}} =  \delta^{(d)}(\vb{r})
\end{equation}
We can use Gauss' theorem to evaluate this.
\begin{equation}
  1 =   \int \dd^d x \ \nabla^2 C_d(\vb{r}) =  \int \dd^d x \ \grad \cdot (\grad C_d(\vb{r})) = \oint\limits_{S} \dd{\vb{S}} \cdot \grad C_d(\vb{r})
\end{equation}
Since the source is at the origin, $C_d$ must be spherically symmetric and hence the gradient is purely radial, that is,
\begin{equation*}
    \grad C_d = \dv{C_d}{r}\ \hat{\vb{r}}
\end{equation*}
Substituting this in the equation above and integrating over a hypersphere, we obtain 
\begin{align*}
     \oint\limits_{S} \dd S \hat{\vb{r}} \cdot \dv{C_d}{r} \hat{\vb{r}}  = \dv{C_d}{r}\oint\limits_S  \dd S = 1
\end{align*}
Since $\dv{C_d}{r}$ on the surface of a hypersphere is a constant, as it is dependent only on $r$, so we pull it out of the integral. The remaining integral is just the surface area of the $d$-dimensional sphere given by $A_d = S_d r^{d-1}$ where $S_d$ is the solid angle (or area of a unit sphere in $d$ dimensions), which implies, 
\begin{equation}
    \dv{C_d}{r} = \frac{1}{S_d r^{d-1}}\ \longrightarrow \ C_d =\frac{1}{S_d} \int\limits_{r\tund{0}}^r \dd r' \frac{1}{(r')^{d-1}} = \begin{cases}
        \frac{1}{S_d(d-2)}\qty[\frac{1}{r^{d-2}} - \frac{1}{r\tund{0}^{d-2}}] & d \neq 2 \\
         \frac{1}{S_d} \ln\qty( \frac{r}{r\tund{0}} )& d=2
    \end{cases}
\end{equation}
where we imposed a lower cutoff to $r$, since the Landau-Ginzburg theory is a coarse-grained theory with a characteristic length scale below which correlations do not have any meaning. Since $\SG_{\perp}\propto C_d$ we have the long-range behaviour ($r\to\infty$) of $\SG_{\perp}(r) $ as,
\begin{equation}
    \SG_{\perp}(r) \sim \begin{cases}
        \frac{1}{r\tund{0}^{d-2}} & d>2  \\[0.3cm] 
       r^{2-d}\to \infty & d<2 \\[0.2cm]
       \ln \qty(\sfrac{r}{r\tund{0}}) \to \infty & d=2
    \end{cases}
\end{equation}
For $d>2$, the phase fluctuation is then finite while for $d\leq 2$, they become asymptotically large. As the phases $\theta\in [0,2\pi)$, the divergence implies that long-range order of the phases is completely washed out for $d\leq 2$. This is the consequence of the \textit{Mermin-Wagner Theorem} which states that systems with continuous symmetry cannot have long-range order in $d\leq 2$ at any finite temperature. Thus for such systems, the lower critical dimension is $d_c=2$.\\

Now, let us examine the correlation function of the order parameter itself, that is, $(not completed)$

